%!TEX root = ../main.tex
\newpage
\section*{謝辞}
\addcontentsline{toc}{section}{謝辞}


本研究を遂行するにあたり、多くの方々から多大なるご支援とご助言をいただきました。

京都大学防災研究所水資源環境センターの小林草平准教授には、研究を継続する貴重な機会を与えてくださっただけでなく、フィールドワークや研究方針について、専門的な知見から温かいご指導をいただきました。
先生のおおらかなお人柄とご配慮がなければ、本研究を完遂することも、研究の楽しさを経験することはできませんでした。
突然のお願いにもかかわらず、私を受け入れてくださり、指導をして頂いたことに心より感謝申し上げます。

空間情報学講座の須崎純一教授には、ゼミでの発表機会や計算資源をご提供いただき、御礼申し上げます。
困難な状況下においても、信念を貫き研究成果を上げられたことは、私にとって大きな自信となりました。
石井順恵助教には、折に触れて親身に相談に乗っていただき、精神的な支えとなってくださったことに深く感謝いたします。 
大場哲治教授にも、学部からの2年間、研究室ゼミでの問いかけを通じて、研究テーマをより深く考える助けをいただきました。
Sameh Kantoush教授には、国際的な研究室ゼミに参加させていただき、実用的かつ示唆に富むアドバイスをいただいたことに感謝いたします。 

幼馴染であり、河川防災システム研究領域研究所の藤井天真さんには、本研究のきっかけをいただいただけでなく、高い志を持って主体的に研究に励む姿から、多大な刺激と勇気をもらいました。 
同期の石黒龍之介さん、岩城直矢さん、風間小次郎さん、濱田邦亮さんをはじめとする研究室の皆様には、日々の研究活動だけでなく、海外活動などを通じて多くの思い出を共有させていただきました。
また、事務手続き面で多大なるご尽力をいただいた佐藤久美子秘書に、感謝申し上げます。

ここに記し切れない多くの方々のご支援とご厚意に支えられ、研究を続けることができました。
% インターンでお世話になったSSSの松尾さんやチームの方々には研究の基礎と本質的に考える思考力を教えていただきました。
% 高校時代の友人は、日常のみならず、研究に関する刺激も与えてくれました。
% 重慶サマーキャンプや、ACRSでのサマースクールで出会った世界中の学生とは、刺激的で楽しい時間を共有できました。
% 就職活動を通して出会った仲間にも、様々な刺激をもらいました。
最後に、大学院進学を支え、今日まで温かく見守ってくれた両親に、感謝を捧げます。
